%% sections/conclusion.tex
\section{Conclusion}
\label{sec:conclusion}

We replicate the DLM~2005 bubble experiment computationally,
confirming that a procedural experience counter---not a cognitive
model of learning---is sufficient to kill the bubble.  Embedding
Lopez-Lira utility agents into the same market produces similar
attenuation under an experience-weighted algorithmic blend.
Replacing the blend with a language-model call yields one clean
finding: the LLM's internal representation of risk preferences is
already robust to whether the closed-form utility expression appears
in the prompt.  Plans~II and~III are indistinguishable in the
standard bubble statistics.  The downstream expected-utility score
absorbs most upstream belief variation, so market dynamics depend more
on the scoring stage and order-book mechanics than on the specific
belief channel.

Three limitations bound the results.  The qualitative ranges in
Table~\ref{tab:mktq} require large-$N$ replication with formal
permutation tests.  The LLM update fires once per period, capping
within-period effects.  Different model families may encode risk
differently; cross-model comparisons are needed.

Several extensions are natural.  A Plan~IV using trust-weighted (rather
than equal-weighted) peer claims would test whether the trust matrix
can substitute for the LLM channel.  Running Plans~II and~III across
model families would establish whether the near-equivalence is
model-general.  A factorial design crossing deception modes with risk
shares would decompose their joint effect on bubble formation.
